% nodos distinguidos
El objetivo de esta sección es analizar las direcciones IP provenientes de los paquetes ARP de las tres redes. En particular, buscamos distinguir símbolos y hosts de las fuentes creadas según la cantidad de información que proveen. Para ello, capturamos hasta $1000$ tramas por red bajo un contexto de uso normal de los dispositivos, en diferentes horarios del día. 

En un primer lugar, observamos que en todas las fuentes los símbolos incluyendo las IPs propias del router y las máquinas utilizadas eran los más probables: en promedio, representaban un $70.6\%$ de los paquetes a lo largo de las redes, conformando un $90.2\%$ de los paquetes en las redes WiFi y un $34.1\%$ en la red Ethernet. Creemos que esta diferencia puede deberse a las particularidades de las redes¡. Mientras que en ambas redes WiFi habían pocos hosts además de los medidores, existían más de 10 hosts conectados al router de la red Ethernet, tanto por WiFi como por cable. De este modo, de requerir actualizar la tabla de direcciones el router debía de enviar los mensajes ARP a ambas redes y, en consecuente, disminuyó la probabilidad de aparición de los símbolos mencionados.
